% NOTES ONLY, to be written. Trimmed 2026-11-07 from thesis/methodology/05a_curvature.tex
% (full draft in git, commit e4eb857). Implementation: CGPS_data/thesis_final/turning_feature.py
% (mean_turning, segment_turning) on each artery from split_trunks.py; figure from
% turning_feature_schematic.py.
%
% The subsection should contain:
% - Curvature as mean turning angle per mm, a form of average curvature, the tortuosity measure most
%   closely, although weakly, related to low WSS (kashyap2022tortuosity; simulation study, keep
%   "although weakly").
% - Turning angle at each point between two 5 mm chords, chord ends by linear interpolation;
%   atan2 form (eq:method-turn), unsigned in [0, pi].
% - Noise: a positional error sigma gives spurious turning that falls with the square of the chord
%   length; 5 mm chord as in vanzandwijk2019curvature. No further smoothing; no values within 5 mm of
%   either artery end. TODO: confirm with Kit that the tracer centerlines arrive smoothed.
% - Profile kappa = theta / l in rad/mm (eq:method-profile, referred to by curvature_fpca), computed
%   on the whole artery and then divided at the segment boundaries.
% - Segment value T_s = arc-length mean of the profile (eq:method-tl), trapezoid rule; fewer than two
%   profile points -> missing.
% - Difference from Tello Ayala et al. (2D, consecutive points): 3D, fixed chord, normalised by
%   length, so independent of point spacing.
% - Figure fig:method-turning.
\subsection{Turning-angle curvature}
\label{sec:method-curvature}

\todonote{To be written: turning angle over 5\,mm chords, the curvature profile and its arc-length mean per segment.}
